\documentclass{beamer}

% Theme choice
\usetheme{Madrid} % You can change to e.g., Warsaw, Berlin, CambridgeUS, etc.

% Encoding and font
\usepackage[utf8]{inputenc}
\usepackage[T1]{fontenc}

% Graphics and tables
\usepackage{graphicx}
\usepackage{booktabs}

% Code listings
\usepackage{listings}
\lstset{
basicstyle=\ttfamily\small,
keywordstyle=\color{blue},
commentstyle=\color{gray},
stringstyle=\color{red},
breaklines=true,
frame=single
}

% Math packages
\usepackage{amsmath}
\usepackage{amssymb}

% Colors
\usepackage{xcolor}

% TikZ and PGFPlots
\usepackage{tikz}
\usepackage{pgfplots}
\pgfplotsset{compat=1.18}
\usetikzlibrary{positioning}

% Hyperlinks
\usepackage{hyperref}

% Title information
\title{Chapter 9: Data Structures II: Dictionaries}
\author{Teaching Assistant}
\institute{University Name}
\date{\today}

\begin{document}

\frame{\titlepage}

\begin{frame}[fragile]
    \frametitle{Chapter 9: Data Structures II: Dictionaries}
    
    \begin{center}
        \huge{\textbf{Chapter 9: Data Structures II: Dictionaries}}
        \vspace{1em}
        
        \large{\textit{Beyond the Index: Modeling Data with Key-Value Pairs}}
    \end{center}
    
    \vfill
    
    \begin{block}{The Core Question}
        We've used lists to store ordered data, accessed by a numeric index (0, 1, 2, ...).
        \vspace{0.5em}
        
        \textbf{Problem:} How do we efficiently find a student's grade by their \textit{name} instead of their position in a list?
    \end{block}
    
\end{frame}

\begin{frame}[fragile]
    \frametitle{The Challenge: Finding Data Efficiently}

    \begin{columns}[T]
        \begin{column}{0.5\textwidth}
            \begin{block}{1. The "Old Way" with Lists}
                We must iterate through the entire list to find a match. This is slow and cumbersome.
                
                \begin{lstlisting}[language=Python]
% The "old way" with lists
student_grades = [['Alice', 95], 
                  ['Bob', 87], 
                  ['Charlie', 92]]

# To find Charlie's grade, we must search:
charlie_grade = 0
for student in student_grades:
    if student[0] == 'Charlie':
        charlie_grade = student[1]
        break
% Inefficient for large datasets!
                \end{lstlisting}
            \end{block}
        \end{column}
        
        \begin{column}{0.5\textwidth}
            \begin{block}{2. The Solution: Dictionaries!}
                A dictionary stores data in \textbf{key-value pairs}. We use a unique key for an instant lookup.
                
                \begin{lstlisting}[language=Python]
% The "new way" with dictionaries
grades_dict = {'Alice': 95, 
               'Bob': 87, 
               'Charlie': 92}

# Instant lookup using the key:
charlie_grade = grades_dict['Charlie']
# Returns 92

% Fast, readable, and efficient!
                \end{lstlisting}
            \end{block}
        \end{column}
    \end{columns}
\end{frame}

\begin{frame}[fragile]
    \frametitle{What You'll Learn in This Chapter}

    \begin{center}
        \large\textbf{Our Roadmap for Mastering Dictionaries}
    \end{center}
    \vfill

    \begin{itemize}
        \item \textbf{Core Concepts:} Understand the key-value model and when to choose a dictionary over a list.
        \vspace{1em}
        
        \item \textbf{Fundamental Operations:} Learn how to create, access, add, update, and remove items from dictionaries.
        \vspace{1em}
        
        \item \textbf{Iteration \& Methods:} Discover powerful techniques for looping through dictionaries and using their built-in methods.
        \vspace{1em}
        
        \item \textbf{Practical Applications:} Model real-world data like JSON, user profiles, and application settings.
    \end{itemize}
    
    \vfill
\end{frame}

\begin{frame}[fragile]
    \frametitle{What is a Dictionary? The Key-Value Pair}
    \begin{block}{The Core Concept: Mappings, Not Sequences}
        A \textbf{dictionary} is a fundamental data structure that stores a collection of items.
        \begin{itemize}
            \item Unlike a list, which is an \textit{ordered sequence} accessed by a numerical index (0, 1, 2, ...), a dictionary is an \textit{unordered collection}.
            \item It stores data as \textbf{key-value pairs}.
        \end{itemize}
    \end{block}

    \begin{alertblock}{Key and Value}
        \begin{itemize}
            \item \textbf{Key:} A unique identifier used to look up data.
            \item \textbf{Value:} The data associated with a key.
        \end{itemize}
        The primary purpose is to retrieve a value quickly when you know its key.
    \end{alertblock}

    \begin{exampleblock}{The Dictionary Analogy}
        Think of a real dictionary: you use a \textbf{word (the key)} to instantly find its \textbf{definition (the value)}. You don't read the whole book to find it!
    \end{exampleblock}
\end{frame}

\begin{frame}[fragile]
    \frametitle{What is a Dictionary? -- Key Rules \& Characteristics}
    Dictionaries have a few simple but important rules:

    \begin{enumerate}
        \item \textbf{Keys Must Be Unique}
        \begin{itemize}
            \item Every key in a dictionary must be unique. You cannot have two identical keys.
            \item If you add a key that already exists, you will simply \textbf{update} the value associated with it.
            \item \textit{Analogy: A phonebook doesn't have two separate entries for the exact same person.}
        \end{itemize}
        \vspace{0.5cm}

        \item \textbf{Values Can Be Duplicated}
        \begin{itemize}
            \item You can have multiple different keys that point to the same value.
            \item \textit{Example: Two different student IDs (keys) could both map to the name "Alex" (value).}
        \end{itemize}
    \end{enumerate}

    \begin{block}{Fast and Efficient Lookups}
        Dictionaries are highly optimized for data retrieval. Accessing a value using its key is extremely fast, regardless of the dictionary's size.
    \end{block}
\end{frame}

\begin{frame}[fragile]
    \frametitle{What is a Dictionary? -- Real-World Analogies}
    Here are other ways to visualize key-value pairs in action:

    \begin{block}{Phonebook}
        \begin{itemize}
            \item \textbf{Key}: Person's Name
            \item \textbf{Value}: Phone Number
        \end{itemize}
    \end{block}

    \begin{block}{Student Records}
        \begin{itemize}
            \item \textbf{Key}: Student ID (e.g., \texttt{900123456})
            \item \textbf{Value}: All Student Information (Name, Major, GPA, etc.)
        \end{itemize}
    \end{block}

    \begin{block}{Website User Profile}
        A dictionary is perfect for storing structured data like a user profile:
        \begin{itemize}
            \item \textbf{Key}: \texttt{'username'} $\rightarrow$ \textbf{Value}: \texttt{'alex\_c'}
            \item \textbf{Key}: \texttt{'user\_id'} $\rightarrow$ \textbf{Value}: \texttt{101}
            \item \textbf{Key}: \texttt{'is\_active'} $\rightarrow$ \textbf{Value}: \texttt{True}
        \end{itemize}
    \end{block}
\end{frame}

\begin{frame}[fragile]
    \frametitle{Syntax: Creating a Dictionary - The Basics}
    Dictionaries are created using curly braces \texttt{\{\}} with key-value pairs separated by colons \texttt{:}.

    \begin{block}{Core Syntax Elements}
        \begin{itemize}
            \item \textbf{Curly Braces \texttt{\{\}}}: Mark the beginning and end of the dictionary.
            \item \textbf{Colons \texttt{:}}: Separate a key from its value.
            \item \textbf{Commas \texttt{,}}: Separate individual key-value pairs.
        \end{itemize}
    \end{block}

    \vfill

    \textbf{General Structure:}
    \begin{verbatim}
{ 'key1': value1, 'key2': value2, ... }
    \end{verbatim}

    \textbf{Creating an Empty Dictionary:}
    \begin{verbatim}
empty_dict = {}
    \end{verbatim}
\end{frame}

\begin{frame}[fragile]
    \frametitle{Syntax: Creating a Dictionary - Code Example}
    Let's create a dictionary to represent a user profile, storing different types of information.

    \begin{block}{Example: \texttt{user\_profile}}
\begin{verbatim}
# A dictionary mapping string keys to various values
user_profile = {
    'username': 'alex_c',
    'user_id': 101,
    'is_active': True,
    'topics': ['python', 'data', 'design']
}
\end{verbatim}
    \end{block}

    \begin{itemize}
        \item \texttt{'username': 'alex\_c'} \hfill (String key, String value)
        \item \texttt{'user\_id': 101} \hfill (String key, Integer value)
        \item \texttt{'is\_active': True} \hfill (String key, Boolean value)
        \item \texttt{'topics': [...]} \hfill (String key, List value)
    \end{itemize}
\end{frame}

\begin{frame}[fragile]
    \frametitle{Syntax: Creating a Dictionary - Key Rules}
    There are two fundamental rules regarding keys and values.

    \begin{block}{Rule 1: Keys Must Be Immutable}
        Keys must be of a data type that cannot be changed. This ensures that a key remains a constant reference.
        \begin{itemize}
            \item \textbf{Allowed Key Types:} Strings, Numbers, Tuples.
            \item \textbf{Not Allowed:} Lists, other Dictionaries (because they are mutable).
        \end{itemize}
    \end{block}

    \begin{block}{Rule 2: Values Can Be Anything}
        A value can be of any data type, which makes dictionaries extremely flexible.
        \begin{itemize}
            \item You can store strings, numbers, booleans, lists, tuples, and even nest other dictionaries as values.
        \end{itemize}
    \end{block}

    \begin{alertblock}{Best Practice: Readability}
        For dictionaries with multiple pairs, writing each key-value pair on a new line (as in the \texttt{user\_profile} example) greatly improves code readability.
    \end{alertblock}
\end{frame}

\begin{frame}[fragile]
    \frametitle{Accessing Values - Part 1}

    You can retrieve a value from a dictionary by providing its key inside square brackets \texttt{[]}. This is similar to indexing a list, but you use a key instead of a numerical index.

    \begin{itemize}
        \item \textbf{Syntax:} \texttt{dictionary\_name[key]}
        \item \textbf{Analogy:}
        \begin{itemize}
            \item List: \texttt{grades[0]} $\rightarrow$ "What is at position 0?"
            \item Dictionary: \texttt{user\_profile['username']} $\rightarrow$ "What is the value for the key 'username'?"
        \end{itemize}
    \end{itemize}

    \textbf{Code Example:}
    \begin{lstlisting}[language=Python]
# Assume user_profile is defined:
user_profile = {
    'username': 'alex_c',
    'user_id': 101,
    'is_active': True,
    'topics': ['python', 'data', 'design']
}

# Get the value associated with the key 'username'
name = user_profile['username']
print(f"Username: {name}")

# Get the list of topics
user_topics = user_profile['topics']
print(f"Topics: {user_topics}")
    \end{lstlisting}
    \textbf{Output:}
    \begin{verbatim}
Username: alex_c
Topics: ['python', 'data', 'design']
    \end{verbatim}
\end{frame}

\begin{frame}[fragile]
    \frametitle{Accessing Values - Part 2: The KeyError Exception}

    One of the most common errors when working with dictionaries is the \texttt{KeyError}.

    \begin{block}{What is a KeyError?}
        Python raises a \texttt{KeyError} when you try to access a key that does \textbf{not} exist in the dictionary. It's Python's way of saying, "I looked for that key, but I couldn't find it!"
        \vspace{0.5em}
        
        This is a helpful feature, as it prevents your program from silently failing or using incorrect data.
    \end{block}

    \textbf{Example of what NOT to do:}
    The \texttt{user\_profile} dictionary does not have an \texttt{'email'} key. Attempting to access it will crash the program.

    \begin{lstlisting}[language=Python]
# This line will raise a KeyError!
print(user_profile['email'])
    \end{lstlisting}

    \textbf{Resulting Traceback:}
    \begin{verbatim}
Traceback (most recent call last):
  File "<stdin>", line 1, in <module>
KeyError: 'email'
    \end{verbatim}
\end{frame}

\begin{frame}[fragile]
    \frametitle{Adding and Updating Items}
    \subtitle{The "Upsert" Operation}
    
    You can add new key-value pairs or update existing ones using the same square bracket assignment syntax.
    
    \begin{block}{The Core Principle: One Syntax for Two Jobs}
        The assignment syntax \texttt{dictionary[key] = value} automatically handles both adding and updating.
    \end{block}
    
    \vfill
    
    \textbf{How Python Decides:}
    \begin{itemize}
        \item \textbf{If the key exists:} The existing value is replaced.
        \begin{itemize}
            \item This is an \textbf{Update}.
        \end{itemize}
        \item \textbf{If the key does not exist:} A new key-value pair is created.
        \begin{itemize}
            \item This is an \textbf{Add} (or Insert).
        \end{itemize}
    \end{itemize}
    
\end{frame}

\begin{frame}[fragile]
    \frametitle{Adding and Updating Items - Example}

    \textbf{Code Example:}
    \begin{verbatim}
# Assume this is our initial dictionary
user_profile = {
    'username': 'alex_c', 'user_id': 101, 
    'is_active': True, 'topics': ['Python', 'Data Science']
}

# Add a new key-value pair
# The key 'email' does not exist, so it is created.
user_profile['email'] = 'alex@example.com'

# Update an existing value
# The key 'is_active' exists, so its value is replaced.
user_profile['is_active'] = False

print(user_profile)
# {'username': 'alex_c', 'user_id': 101, 'is_active': False, 
# 'topics': [...], 'email': 'alex@example.com'}
    \end{verbatim}
    
    \begin{block}{Key Takeaways}
        \begin{itemize}
            \item \textbf{Dictionaries are Mutable}: You can freely add, update, and remove items after creation.
            \item \textbf{One Syntax for Two Jobs}: The \texttt{dict[key] = value} assignment is an all-in-one tool.
            \item \textbf{No \texttt{KeyError} on Assignment}: Unlike accessing, assigning to a non-existent key is a valid operation that creates it.
        \end{itemize}
    \end{block}

\end{frame}

\begin{frame}[fragile]
    \frametitle{Removing Items}
    \framesubtitle{Using the \texttt{del} Keyword}

    To remove a key-value pair from a dictionary, you can use the \texttt{del} keyword. This action is permanent and modifies the dictionary in-place.

    \begin{itemize}
        \item \textbf{Purpose:} To remove a specific entry from a dictionary using its key.
        \item \textbf{Mechanism:} \texttt{del} is a general Python statement that deletes objects.
        \item \textbf{Syntax:} \texttt{del dictionary\_name['key\_to\_remove']}
    \end{itemize}

    \begin{block}{Code Example}
\begin{verbatim}
user_profile = {
    'username': 'alex_c', 'user_id': 101, 
    'is_active': False, 
    'topics': ['python', 'data science', 'ai']
}

# Let's remove the user_id
del user_profile['user_id']

print(user_profile)
# {'username': 'alex_c', 'is_active': False, 
#  'topics': ['python', 'data science', 'ai']}
\end{verbatim}
    \end{block}
\end{frame}

\begin{frame}[fragile]
    \frametitle{Removing Items - Errors \& Alternatives}

    \begin{block}{Caution: The \texttt{KeyError}}
        Attempting to \texttt{del} a key that does not exist will raise a \texttt{KeyError}, crashing your program if not handled.
\begin{verbatim}
# The key 'last_login' does not exist
del user_profile['last_login']

# Raises: KeyError: 'last_login'
\end{verbatim}
        \textbf{Key Takeaway:} Always ensure a key exists before trying to delete it (e.g., \texttt{if 'key' in my\_dict:}).
    \end{block}

    \begin{block}{Alternative: The \texttt{.pop()} Method}
        The \texttt{.pop()} method is another way to remove an item. Its key advantage is that it \textbf{removes the item and returns its value}.
\begin{verbatim}
# Remove email and store it in a variable
removed_email = user_profile.pop('email')

print(f"Removed email: {removed_email}")
# > Removed email: alex@example.com

print(user_profile)
# {'username': 'alex_c', 'is_active': False, 
#  'topics': [...]}
\end{verbatim}
        \texttt{.pop()} can also take a default value (e.g., \texttt{.pop('key', None)}) to avoid a \texttt{KeyError}, making it a safer choice.
    \end{block}
\end{frame}

\begin{frame}[fragile]
    \frametitle{Iterating Through a Dictionary}
    \framesubtitle{Method 1: Iterating Through Keys (Default)}

    Just like lists, you can use a \texttt{for} loop to process the contents of a dictionary. By default, a \texttt{for} loop over a dictionary iterates through its keys.

    \begin{block}{Example: Looping Through Keys}
\begin{lstlisting}[language=Python]
user_profile = {
    'username': 'alex_c', 
    'is_active': True, 
    'topics': ['Python', 'AI', 'Data Science']
}

# You can use the key to access its value inside the loop
print("--- Profile Details ---")
for key in user_profile:
    value = user_profile[key]
    print(f"The value for '{key}' is: {value}")
\end{lstlisting}
    \end{block}
    
    \begin{block}{Output}
\begin{verbatim}
--- Profile Details ---
The value for 'username' is: alex_c
The value for 'is_active' is: True
The value for 'topics' is: ['Python', 'AI', 'Data Science']
\end{verbatim}
    \end{block}
\end{frame}

\begin{frame}[fragile]
    \frametitle{Iterating Through a Dictionary}
    \framesubtitle{Method 2: Key-Value Pairs with \texttt{.items()}}

    The most useful and expressive way to loop is with the \texttt{.items()} method. It returns key-value pairs that can be "unpacked" into two variables directly in the \texttt{for} loop. This is the most common pattern.

    \begin{block}{Code: Using \texttt{.items()}}
\begin{lstlisting}[language=Python]
# The .items() method is efficient and clear
print("--- Full Profile using .items() ---")
for key, value in user_profile.items():
    print(f"{key.capitalize()}: {value}")
\end{lstlisting}
    \end{block}

    \begin{block}{Output}
\begin{verbatim}
--- Full Profile using .items() ---
Username: alex_c
Is_active: True
Topics: ['Python', 'AI', 'Data Science']
\end{verbatim}
    \end{block}
\end{frame}

\begin{frame}[fragile]
    \frametitle{Dictionary Iteration: Summary}

    Here are the three primary ways to iterate through a dictionary:

    \begin{itemize}
        \item \textbf{Loop over Keys (Default):} Use when you only need the keys.
        \begin{lstlisting}[language=Python, basicstyle=\small\ttfamily]
for key in my_dict:
    # ... use key
        \end{lstlisting}

        \item \textbf{Loop over Values:} Use when you only need the values, not the keys.
        \begin{lstlisting}[language=Python, basicstyle=\small\ttfamily]
for value in my_dict.values():
    # ... use value
        \end{lstlisting}

        \item \textbf{Loop over Both (Recommended):} The most powerful and readable pattern.
        \begin{lstlisting}[language=Python, basicstyle=\small\ttfamily]
for key, value in my_dict.items():
    # ... use both key and value
        \end{lstlisting}
    \end{itemize}

    \begin{block}{Key Takeaway}
    Always use \texttt{my\_dict.items()} when you need both the key and the value. It is more efficient and makes your code easier to understand.
    \end{block}
\end{frame}

\begin{frame}[fragile]
    \frametitle{Checking for Keys - The \texttt{KeyError} Trap}

    When working with dictionaries, one of the most common errors is the \texttt{KeyError}.
    \begin{itemize}
        \item It occurs when you try to access a key that \textbf{does not exist}.
        \item This error will immediately \textbf{halt the execution} of your program.
    \end{itemize}

    \begin{block}{Example of what causes a \texttt{KeyError}}
        \begin{lstlisting}[language=Python]
% Define a simple user profile
user_profile = {'username': 'jdoe', 'id': 101}

# Trying to access the 'email' key, which doesn't exist.
# This line will immediately cause a KeyError!
print(user_profile['email'])
        \end{lstlisting}
    \end{block}
    \vspace{1em}
    \textbf{Goal:} Write robust code that anticipates missing keys and handles them gracefully.
\end{frame}

\begin{frame}[fragile]
    \frametitle{Checking for Keys - The \texttt{in} Keyword Solution}
    To avoid a \texttt{KeyError}, we can practice "defensive programming" by first checking if a key exists using the \texttt{in} keyword.

    \begin{itemize}
        \item The expression \texttt{key in dictionary} returns a boolean:
        \begin{itemize}
            \item \texttt{True} if the key is present.
            \item \texttt{False} if the key is not present.
        \end{itemize}
        \item You can also use \texttt{key not in dictionary} to check for a key's absence.
        \item \textbf{Important:} The \texttt{in} operator checks for membership in the dictionary's \textbf{keys}, not its values.
    \end{itemize}

    \begin{block}{Code Example: Safe Key Access}
        \begin{lstlisting}[language=Python]
if 'email' in user_profile:
    print(f"Email found: {user_profile['email']}")
else:
    print("Email key does not exist.")

# Check for a non-existent key
if 'last_login' not in user_profile:
    print("User has not logged in yet.")
        \end{lstlisting}
    \end{block}
\end{frame}

\begin{frame}[fragile]
\frametitle{When to Use a Dictionary vs. a List}
\framesubtitle{Choosing the Right Tool for the Job}

Choosing the right data structure is crucial for writing efficient and readable code. Your choice depends on how you need to store and access your data.

\vspace{0.5em}

\begin{columns}[T] % The 'T' option aligns the tops of the columns

%=============== COLUMN 1: LIST ===============
\begin{column}{0.5\textwidth}
    \begin{block}{Use a List when...}
    \begin{itemize}
        \item \textbf{Order Matters}: You have a collection of items where the sequence is important.
        \item \textbf{Access by Position}: You need to get elements by their numerical index (0, 1, 2, ...).
        \item \textbf{Mental Model}: Think of a numbered list, a queue, or a sequence of steps.
    \end{itemize}
    \end{block}
    
    \begin{block}{Example: Steps in a Recipe}
\begin{verbatim}
recipe_steps = [
    "Preheat oven to 350 F",
    "Mix dry ingredients",
    "Combine and pour into pan"
]

# Get the second step (index 1)
second_step = recipe_steps[1]
print(second_step)
# Output: Mix dry ingredients
\end{verbatim}
    \end{block}
\end{column}

%=============== COLUMN 2: DICTIONARY ===============
\begin{column}{0.5\textwidth}
    \begin{block}{Use a Dictionary when...}
    \begin{itemize}
        \item \textbf{Labels Matter}: You need to store data as key-value pairs, where each item has a unique name (key).
        \item \textbf{Access by Identifier}: You need to look up a value based on its key, not its position.
        \item \textbf{Mental Model}: Think of a phone book, a real dictionary, or a user profile.
    \end{itemize}
    \end{block}
    
    \begin{block}{Example: Application Settings}
\begin{verbatim}
app_settings = {
    "theme": "dark",
    "font_size": 14,
    "notifications_enabled": True
}

# Get the theme by its key
current_theme = app_settings["theme"]
print(current_theme)
# Output: dark
\end{verbatim}
    \end{block}
\end{column}

\end{columns}

\vspace{1em}

\begin{alertblock}{Key Question to Ask Yourself}
"Do I need to know the \textbf{order} of my items (use a \textbf{List}), or do I need to look up items by a specific \textbf{name} or \textbf{ID} (use a \textbf{Dictionary})?"
\end{alertblock}

\end{frame}

\begin{frame}[fragile]
    \frametitle{Chapter Summary - Core Concepts}

    \begin{block}{1. The Key-Value Pair}
        The fundamental concept that defines a dictionary. It's an \textbf{unordered} collection where each value is associated with a unique identifier (the key).
        \begin{itemize}
            \item \textbf{Keys are Unique:} No two keys can be the same in a single dictionary.
            \item \textbf{Keys are Immutable:} Must be a data type that cannot change (e.g., string, number, tuple).
        \end{itemize}
    \end{block}

    \begin{block}{2. Creation Syntax: \texttt{\{\}}}
        We use curly braces to create dictionaries and a colon to separate a key from its value.
        \begin{lstlisting}[language=Python]
% A dictionary representing a student record
student = {
    'student_id': 93451,
    'name': 'Maria Garcia',
    'major': 'Physics',
    'gpa': 3.7
}
        \end{lstlisting}
    \end{block}
\end{frame}

\begin{frame}[fragile]
    \frametitle{Chapter Summary - Operations \& Use Cases}
    \begin{block}{3. Core Operations}
        All major operations use the key to identify the item.
        \begin{columns}[T]
            \begin{column}{0.4\textwidth}
                \begin{itemize}
                    \item \textbf{Access}: \texttt{student['name']}
                    \item \textbf{Add / Update}: \\ \texttt{student['gpa'] = 3.8}
                    \item \textbf{Delete}: \texttt{del student['major']}
                \end{itemize}
            \end{column}
            \begin{column}{0.6\textwidth}
                \begin{lstlisting}[language=Python]
# Accessing a value
student_name = student['name']

# Updating a value
student['gpa'] = 3.8

# Adding a new pair
student['grad_year'] = 2025

# Deleting a pair
del student['major']
                \end{lstlisting}
            \end{column}
        \end{columns}
    \end{block}

    \begin{block}{4. When to Use Dictionaries}
        Use a dictionary whenever your data has inherent labels or identifiers you need to look up.
        \begin{itemize}
            \item \textbf{User Profiles:} Look up user data by \texttt{'username'} or \texttt{'user\_id'}.
            \item \textbf{Configuration Settings:} Store and retrieve settings by name, e.g., \texttt{'theme'}.
            \item \textbf{Records:} Model any real-world object like a product, car, or database entry.
        \end{itemize}
    \end{block}
\end{frame}


\end{document}